%% kapa: exam=1 round=34 year=2023 period=1 form=A subject=economics subject_ko=경제학원론 seq=35 number=75
%% source: Q-Net 감정평가사 기출문제 (cst00309) artlSeq=5211789 file=raw/1008/5211789/2023년도 제34회 감정평가사 1차시험 1교시.pdf
%% conversion: kordoc 4.24.51 -> tools/gen1.py (auto); review=proofread-sonnet; flags=none
\begin{question}{75}
A국의 중앙은행은 필립스곡선, 성장률과 실업률의 관계, 이자율 준칙에 따라 이자율을 결정한다. 현재 목표물가상승률이 2\%, 자연실업률이 3\%이고, 국내총생산은 잠재국내총생산, 물가상승률은 목표물가상승률, 그리고 실업률은 자연실업률과 같다고 가정할 때, 이에 관한 설명으로 옳지 않은 것은? \par
(단, $r$, $\pi$, $\pi^e$, $\pi^T$, $u$, $u_n$, $u_{-1}$, $Y$, $Y^P$는 각각 이자율, 물가상승률, 기대물가상승률, 목표물가상승률, 실업률, 자연실업률, 전기의 실업률, 국내총생산, 잠재국내총생산이다.)
\begin{framedtext}
○ 필립스곡선: $\pi = \pi^e - 0.5(u - u_n)$ \par
○ 이자율 준칙: $r = \pi + 2.0\% + 0.5(\pi - \pi^T) + 0.5G$ \par
○ 성장률과 실업률의 관계: 국내총생산의 성장률 = $3\% - 2(u - u_{-1})$ \par
○ $G = \dfrac{(Y - Y^P)}{Y^P} \times 100$ \par
\end{framedtext}
\begin{choices}
  \item 현재 이자율은 4\%이다.
  \item 현재 기대물가상승률은 2\%이다.
  \item 실업률이 5\%로 상승하고 기대물가상승률이 변화하지 않았다면, 물가상승률은 1\%이다.
  \item 기대물가상승률이 3\%로 상승하면, 이자율은 5.5\%이다.
  \item 실업률이 1\%로 하락하고, 기대물가상승률이 3\%로 상승하면, 이자율은 7\%이다.
\end{choices}
\end{question}
\answer{5}
